% Л12 -- Кэширование, логирование и наблюдаемость. Slides.
%
% Companion to the spoken script in ../L12-caching-and-observability.md, which
% is the source of truth: the slides carry diagrams, code, real terminal output
% and short marks, and the delivery lives in the script only -- this deck
% carries no speaker notes. Slide text is in Russian because that is the
% content; everything else in this file is English.
%
% Build:  make          (latexmk -pdf; engine: pdflatex)
%
% The C# listings were compiled on dotnet SDK 10.0 in a scratch
% `dotnet new web -n HelloWorld` (the AppHost listing in a scratch Aspire
% solution); terminal output on the slides is captured from runs of that
% project. See ../CLAUDE.md for the verification table.
%
% House rules: ../CLAUDE.md and ../../../CLAUDE.md. Theme: ../../../../theme/.

\documentclass[aspectratio=169,10pt,t]{beamer}

\usetheme{dotnet}
\usetikzlibrary{arrows.meta,positioning,calc,decorations.pathreplacing,fit}
\usepackage{array}   % ragged-right table columns

% A ragged-right paragraph column: justified cells in narrow columns hyphenate
% badly, and a split word is the one thing the back row cannot read.
\newcolumntype{L}[1]{>{\raggedright\arraybackslash}p{#1}}

\dotnetlecture{Л12}

\title[Л12. Кэширование, логирование и наблюдаемость]{Кэширование, логирование
  и наблюдаемость}
\author{Дмитрий Курочкин}
\date{}

% Code sizes, one step down from \footnotesize. Smaller than 6.6pt does not
% carry to the back of the room.
\newcommand{\codesm}{\ttfamily\fontsize{7.6}{9.0}\selectfont}
\newcommand{\codexs}{\ttfamily\fontsize{7.0}{8.3}\selectfont}
\newcommand{\codexxs}{\ttfamily\fontsize{6.8}{7.8}\selectfont}
\newcommand{\termsm}{\ttfamily\fontsize{6.6}{7.9}\selectfont}

% A mark at the start of a listing line, for labels set beside the listing.
% Used through the listing's escapechar; the labels are drawn by an overlay
% picture after the listing, so the build needs two passes (latexmk does that).
\newcommand{\mk}[1]{\tikz[remember picture,overlay]\coordinate (#1) at (0,0.55ex);}

\tikzset{
  rawbrace/.style={decorate,decoration={brace,amplitude=4pt},hairline,
                   line width=0.8pt},
  rawlab/.style={anchor=west,font=\footnotesize,inner sep=0pt},
  box/.style={draw=hairline,line width=0.8pt,fill=codebg,inner sep=4pt,
              align=center,minimum height=0.9cm},
  arr/.style={-{Straight Barb[length=1.6mm]},muted,line width=0.7pt},
  arra/.style={-{Straight Barb[length=1.8mm]},accent,line width=0.9pt},
  arrl/.style={-{Straight Barb[length=1.8mm]},alt,line width=0.9pt},
}

\begin{document}
\shorthandoff{"}

% ------------------------------------------------------------------ title --
% 1
\begin{frame}[plain]
  \titlepagednc{Лекция 12}{Кэширование, логирование и наблюдаемость}{}{Дмитрий
    Курочкин}
\end{frame}

% ============================================================ 1. CACHING ==
\section{1. Кэширование}

% 2
\begin{frame}[plain]
  \sectionopener{1}{Кэширование}{От простого к \hi{распределённому}.}
\end{frame}

% 3
\begin{frame}[plain]
  \vote{Что такое кэш?}
\end{frame}

% 4
\begin{frame}[plain]
  \statement{Кэш --- это \hi{согласие} отдавать \lo{устаревшие} данные ради
    скорости. Осознанное.}
\end{frame}

% 5
\begin{frame}{Цена и выгода}
  \slidetop
  \versus{Получаем}{%
    \begin{cbul}{accent}
      \item микросекунды вместо миллисекунд
      \item ни сети, ни соединения из пула
      \item 100\,000 запросов к базе --- \hi{один} в минуту
      \item разница в порядках
    \end{cbul}}{Платим}{%
    \begin{cbul}{alt}
      \item устаревание
      \item память
      \item второй источник правды
      \item ошибка «иногда, у части пользователей»
    \end{cbul}}

  \slidegap
  \hrulethin{0.42}
  \rulegap
  \aside{Решили вы, понимая цену, --- это кэш. Случилось само --- это баг
    «изменил, а вижу старое».}
\end{frame}

% 6
\begin{frame}{Что кэшируют}
  \slidetop
  \versus{Кэшируют}{%
    \begin{cbul}{accent}
      \item читают часто, меняют редко
      \item дорого вычислять
      \item за чем ходят наружу
    \end{cbul}}{Не кэшируют}{%
    \begin{cbul}{alt}
      \item записи
      \item уникальное для запроса
      \item то, что и так быстро
    \end{cbul}}

  \slidegap
  \hrulethin{0.42}
  \rulegap
  \aside{Сначала измерить: кэшируется \hi{медленное, частое и одинаковое}. И
    измерять после --- долю попаданий: 90\,\% --- база видит десятую часть
    нагрузки, 10\,\% --- кэш надо убирать.}
\end{frame}

% 7
\begin{frame}{Лестница}
  \slidetop
  {\renewcommand{\arraystretch}{1.55}
  \begin{tabular}{@{}L{2.4cm}L{3.3cm}L{2.7cm}L{3.5cm}@{}}
    & \colhead{accent}{где лежит} & \colhead{accent}{кто видит}
      & \colhead{alt}{когда ломается} \\
    \textbf{кэш ответа}   & ответ целиком, на границе & все клиенты
      & ответ зависит от пользователя \\
    \textbf{в памяти}     & куча процесса             & один экземпляр
      & экземпляров два \\
    \textbf{Redis}        & отдельный сервис          & все экземпляры
      & сеть и сериализация \\
    \hi{\textbf{гибридный}} & память поверх Redis     & экземпляр и все
      & локальная копия отстала \\
  \end{tabular}}

  \slidegap
  \aside{Четыре ступени --- четыре ответа на вопрос «где лежит и кто видит».}
\end{frame}

% 8
\begin{frame}[fragile]{Кэш ответа}
  \slidetop
  \claim{Копия уходит в сокет, \hi{не доходя} до обработчика.}

\begin{lstlisting}[style=csharp,basicstyle=\codesm,escapechar=|]
|\mk{o1}|builder.Services.AddOutputCache();
|\mk{o2}|app.UseOutputCache();

app.MapGet("/products", () => catalog.All())
   .CacheOutput(policy => policy
|\mk{o3}|       .Expire(TimeSpan.FromMinutes(5))
|\mk{o4}|       .SetVaryByQuery("page", "sort")
|\mk{o5}|       .Tag("catalog"));
\end{lstlisting}
  \begin{tikzpicture}[remember picture,overlay,every node/.style={rawlab}]
    \draw[rawbrace] ([xshift=8.15cm,yshift=0.35em]o1) -- ([xshift=8.15cm,yshift=-0.3em]o2);
    \node at ([xshift=8.4cm]$(o1)!0.5!(o2)$) {регистрация и middleware};
    \node at ([xshift=8.4cm]o3) {срок};
    \node at ([xshift=8.4cm]o4) {от чего зависит ключ};
    \node[text=accent] at ([xshift=8.4cm]o5) {сброс группой};
  \end{tikzpicture}

  \slidegap
  \aside{Для того, что одинаково для всех. Запрос с заголовком
    \texttt{Authorization} \lo{не кэшируется}: отдать чужой ответ --- утечка.}
\end{frame}

% 9
\begin{frame}{Нулевая ступень: заголовки HTTP}
  \slidetop
  \claim{Кэш у клиента и прокси --- \hi{ни байта} вашей памяти.}

  \begin{center}
    \begin{tikzpicture}[font=\footnotesize]
      \node[box,minimum width=2.6cm,minimum height=2.6cm] (c) at (0,0) {клиент\\или прокси};
      \node[box,minimum width=2.6cm,minimum height=2.6cm] (s) at (9.4,0) {сервис};
      \draw[arr] ([yshift=0.85cm]s.west) -- node[above] {\texttt{200} и тело,
        \texttt{ETag: "v7"}} ([yshift=0.85cm]c.east);
      \draw[arr] (c.east) -- node[above] {\texttt{If-None-Match: "v7"}} (s.west);
      \draw[arra] ([yshift=-0.85cm]s.west) -- node[above,text=accent]
        {\texttt{304 Not Modified}, без тела} ([yshift=-0.85cm]c.east);
    \end{tikzpicture}
  \end{center}

  \slidegap
  \aside{Экономия и на вычислении, и на передаче. Цена --- считать отпечаток:
    нужна версия данных.}
\end{frame}

% 10
\begin{frame}{Кэш в памяти}
  \slidetop
  \versus{Даёт}{%
    \begin{cbul}{accent}
      \item объект в той же куче
      \item ссылка, ноль сериализации
      \item «получить или создать»
      \item срок: абсолютный или скользящий
    \end{cbul}}{Требует}{%
    \begin{cbul}{alt}
      \item предел по числу записей
      \item или короткий срок жизни
      \item приоритет вытеснения
      \item иначе --- утечка с отложенным взрывом
    \end{cbul}}

  \slidegap
  \hrulethin{0.42}
  \rulegap
  \aside{Ключ на каждого пользователя за всё время --- куча растёт до потолка, а
    на графике это выглядит как проблема сборщика мусора.}
\end{frame}

% 11
\begin{frame}[plain]
  \vote{Два экземпляра за балансировщиком. Профиль изменили на первом.
    Следующий запрос попал на второй. Что он увидит?}
\end{frame}

% 12
\begin{frame}{Состояние в памяти экземпляра}
  \slidetop
  \begin{center}
    \begin{tikzpicture}[font=\footnotesize]
      \node[box,minimum width=2.2cm,minimum height=1.0cm] (lb) at (0,0) {балансировщик};
      \node[box,draw=accent,line width=1.0pt,minimum width=4.4cm,minimum height=1.1cm]
        (a) at (5.2,0.95) {\textbf{экземпляр 1}\\кэш: профиль \hi{новый}};
      \node[box,draw=alt,line width=1.0pt,minimum width=4.4cm,minimum height=1.1cm]
        (b) at (5.2,-0.95) {\textbf{экземпляр 2}\\кэш: профиль \lo{старый}, ещё 5 минут};
      \node[box,minimum width=1.8cm,minimum height=1.0cm] (db) at (10.0,0) {база};
      \draw[arra] (lb.east) -- node[above,sloped,text=accent] {изменить} (a.west);
      \draw[arrl] (lb.east) -- node[below,sloped,text=alt] {прочитать} (b.west);
      \draw[arra] (a.east) -- (db.west);
    \end{tikzpicture}
  \end{center}

  \slidegap
  \begin{cbul}{accent}
    \item пользователь видит своё изменение \lo{через раз}
    \item это состояние в памяти --- то, что stateless запрещает
    \item допустимо для того, что не меняется: справочник стран, настройки, шаблоны
  \end{cbul}
\end{frame}

% 13
\begin{frame}{Распределённый кэш: Redis}
  \slidetop
  \claim{Положил один --- прочитал \hi{любой}. Сбросил один --- нет ни у кого.}

  {\renewcommand{\arraystretch}{1.5}
  \begin{tabular}{@{}L{3.6cm}L{9.6cm}@{}}
    \hi{\textbf{значение}}   & байты: объект сериализуется туда и обратно \\
    \hi{\textbf{сеть}}       & не микросекунды, а доли миллисекунды \\
    \hi{\textbf{соединение}} & долгоживущее, одно на приложение --- singleton \\
    \hi{\textbf{роль кэша}}  & без сохранения на диск, вытеснение по памяти \\
  \end{tabular}}

  \slidegap
  \aside{Состояние вынесено из процесса --- сервис снова stateless. Кэш --- то,
    что \lo{можно потерять}.}
\end{frame}

% 14
\begin{frame}{Гибридный кэш}
  \slidetop
  \claim{Локальный кэш в памяти \hi{поверх} распределённого.}

  \begin{center}
    \begin{tikzpicture}[font=\footnotesize,
      tb/.style={box,minimum width=2.75cm,minimum height=1.1cm}]
      \node[tb] (q) at (0,0) {\textbf{запрос}\\по ключу};
      \node[tb,draw=accent,line width=1.0pt] (l1) at (3.55,0) {\textbf{память}\\микросекунды};
      \node[tb] (l2) at (7.1,0) {\textbf{Redis}\\доли миллисекунды};
      \node[tb] (src) at (10.65,0) {\textbf{источник}\\миллисекунды};
      \draw[arr] (q) -- (l1);
      \draw[arr] (l1) -- node[above,font=\scriptsize] {нет} (l2);
      \draw[arr] (l2) -- node[above,font=\scriptsize] {нет} (src);
      \draw[arra] (src.south) -- ++(0,-0.55) -| node[pos=0.25,below,text=accent,
        font=\scriptsize] {положить в оба} (l1.south);
    \end{tikzpicture}
  \end{center}

  \slidegap
  \aside{Горячие ключи читают тысячи раз в секунду --- их выгоднее держать
    локально. Цена --- снова состояние в памяти экземпляра.}
\end{frame}

% 15
\begin{frame}[fragile]{HybridCache}
  \vspace{0.3em}
\begin{lstlisting}[style=csharp,basicstyle=\codexs,escapechar=|]
public sealed class CatalogService(
    HybridCache cache, ICatalogRepository repository)
{
    public ValueTask<IReadOnlyList<ProductDto>> GetForUserAsync(
        string userId, CancellationToken ct) =>
        cache.GetOrCreateAsync(
|\mk{h1}|            key: $"catalog:v2:user:{userId}",
|\mk{h2}|            factory: async token =>
                await repository.LoadForUserAsync(userId, token),
|\mk{h3}|            options: new HybridCacheEntryOptions
                { Expiration = TimeSpan.FromMinutes(5) },
|\mk{h4}|            tags: ["catalog"],
            cancellationToken: ct);

    public ValueTask InvalidateCatalogAsync(CancellationToken ct) =>
|\mk{h5}|        cache.RemoveByTagAsync("catalog", ct);
}
\end{lstlisting}
  \begin{tikzpicture}[remember picture,overlay,every node/.style={rawlab}]
    \node[text=accent] at ([xshift=9.6cm]h1) {ключ с пользователем};
    \node at ([xshift=9.6cm]h2) {фабрика: только если пусто};
    \node at ([xshift=9.6cm]h3) {срок};
    \node at ([xshift=9.6cm]h4) {теги};
    \node[text=accent] at ([xshift=9.6cm]h5) {сброс по тегу};
  \end{tikzpicture}

  \aside{Первый вызов идёт в хранилище, следующие пять минут --- в кэш.}
\end{frame}

% 16
\begin{frame}{Защита от лавины}
  \slidetop
  \claim{Тысяча одновременных промахов --- \hi{один} вызов фабрики.}

  \versus{HybridCache}{%
    \begin{cbul}{accent}
      \item фабрика вызвана один раз
      \item 999 запросов ждут её результат
      \item база видит один запрос
    \end{cbul}}{Наивный кэш}{%
    \begin{cbul}{alt}
      \item фабрика вызвана тысячу раз
      \item тысяча одинаковых запросов
      \item срок истёк --- кэш добивает базу
    \end{cbul}}

  \slidegap
  \hrulethin{0.42}
  \rulegap
  \aside{В Redis лежат байты --- кладите DTO. В памяти лежит сам объект, один на
    всех: \hi{record}, и никаких правок после получения.}
\end{frame}

% 17
\begin{frame}{Ключ по частям}
  \slidetop
  \claim{Всё, от чего зависит ответ, --- \hi{в ключе}.}

  \begin{center}
    \begin{tikzpicture}[font=\footnotesize]
      \node[font=\ttfamily\LARGE,inner sep=0pt,anchor=west] (k1) at (0,0) {catalog};
      \node[font=\ttfamily\LARGE,inner sep=0pt,anchor=west,text=muted] (c1) at (k1.east) {:};
      \node[font=\ttfamily\LARGE,inner sep=0pt,anchor=west] (k2) at (c1.east) {v2};
      \node[font=\ttfamily\LARGE,inner sep=0pt,anchor=west,text=muted] (c2) at (k2.east) {:};
      \node[font=\ttfamily\LARGE,inner sep=0pt,anchor=west,text=accent] (k3) at (c2.east)
        {user:\{userId\}};
      \draw[rawbrace] ([yshift=-0.7em]k1.south east) -- ([yshift=-0.7em]k1.south west);
      \draw[rawbrace] ([yshift=-0.7em]k2.south east) -- ([yshift=-0.7em]k2.south west);
      \draw[rawbrace,accent] ([yshift=-0.7em]k3.south east) -- ([yshift=-0.7em]k3.south west);
      \node[below=1.5em,align=center] at (k1.south) {что это};
      \node[below=1.5em,align=center] at (k2.south) {версия\\формата};
      \node[below=1.5em,align=center,text=accent] at (k3.south) {для кого};
    \end{tikzpicture}
  \end{center}

  \slidegap
  \begin{cbul}{accent}
    \item версия: новый код читает по новым ключам --- выкладка без лавины
    \item имена читаемые: в Redis их будут смотреть глазами
    \item Redis общий --- ключ начинается с имени сервиса
  \end{cbul}
\end{frame}

% 18
\begin{frame}[plain]
  \vote{Ключ \texttt{catalog}, без пользователя. Менеджер видит закупочные цены,
    клиент --- нет. Что произойдёт?}
\end{frame}

% 19
\begin{frame}{Чужой ответ}
  \slidetop
  \claim{Авторизация отработала правильно. Утечку устроил \lo{ключ}.}

  \begin{center}
    \begin{tikzpicture}[font=\footnotesize,
      tb/.style={box,minimum width=3.4cm,minimum height=1.1cm}]
      \node[tb] (m) at (0,0) {\textbf{менеджер}\\запросил первым};
      \node[tb,draw=alt,line width=1.0pt] (k) at (4.8,0) {\textbf{ключ \texttt{catalog}}\\ответ менеджера, 5 минут};
      \node[tb] (u) at (9.6,0) {\textbf{все клиенты}\\видят закупочные цены};
      \draw[arr] (m) -- (k);
      \draw[arrl] (k) -- (u);
    \end{tikzpicture}
  \end{center}

  \slidegap
  \begin{cbul}{accent}
    \item ответ зависит от того, кто спрашивает, --- \hi{он в ключе}, всегда
    \item кэш стоит \hi{после} авторизации, не до неё
  \end{cbul}
\end{frame}

% 20
\begin{frame}{Инвалидация: три стратегии}
  \slidetop
  {\renewcommand{\arraystretch}{1.55}
  \begin{tabular}{@{}L{2.8cm}L{4.8cm}L{5.2cm}@{}}
    & \colhead{accent}{даёт} & \colhead{alt}{цена} \\
    \textbf{по времени} & просто и предсказуемо & устаревает ровно на срок \\
    \textbf{по событию} & изменение видно сразу & знать все места, где данные
      меняются \\
    \textbf{по версии}  & ничего не нужно сбрасывать & версию надо где-то взять;
      старое живёт до срока \\
  \end{tabular}}

  \slidegap
  \aside{Их сочетают: короткий срок как \hi{страховка}, сброс по событию --- где
    изменения известны, версия --- где она и так есть. Идеального ответа нет.}
\end{frame}

% 21
\begin{frame}{Слои, а не альтернативы}
  \slidetop
  \begin{center}
    \begin{tikzpicture}[font=\footnotesize,
      tb/.style={box,minimum width=2.75cm,minimum height=1.15cm}]
      \node[tb] (a) at (0,0)    {\textbf{клиент, прокси}\\заголовки HTTP};
      \node[tb] (b) at (3.55,0) {\textbf{кэш ответа}\\до обработчика};
      \node[tb] (c) at (7.1,0)  {\textbf{гибридный кэш}\\память и Redis};
      \node[tb,draw=accent,line width=1.0pt] (d) at (10.65,0) {\textbf{база}\\малая часть};
      \draw[arr] (a) -- (b);
      \draw[arr] (b) -- (c);
      \draw[arr] (c) -- (d);
      \node[below=0.5em,text=muted,font=\scriptsize] at (a.south) {запрос не пришёл};
      \node[below=0.5em,text=muted,font=\scriptsize] at (b.south) {не дошёл до кода};
      \node[below=0.5em,text=muted,font=\scriptsize] at (c.south) {не дошёл до базы};
    \end{tikzpicture}
  \end{center}

  \slidegap
  \claim{Каждый слой ловит \hi{свою долю} запросов.}

  \aside{Ошибка на любом слое одна и та же --- старое или чужое. И лечится
    одним и тем же: \hi{ключом и сроком}.}
\end{frame}

% 22
\begin{frame}{Границы}
  \slidetop
  {\renewcommand{\arraystretch}{1.5}
  \begin{tabular}{@{}L{4.2cm}L{9.0cm}@{}}
    \hi{\textbf{«не найдено»}}     & да, но секунды: объект может появиться \\
    \lo{\textbf{ошибка}}           & нет: это не результат, а его отсутствие \\
    \lo{\textbf{создание заказа}}  & нет: кэшируются только безопасные чтения \\
    \hi{\textbf{гонка с записью}}  & короткий срок: старое живёт минуту, а не сутки \\
    \hi{\textbf{холодный старт}}   & прогреть до того, как экземпляр скажет «готов» \\
  \end{tabular}}

  \slidegap
  \aside{Идемпотентность и кэшируемость --- два свойства одной оси.}
\end{frame}

% 23
\begin{frame}[plain]
  \keyline{Кэш --- согласие на устаревшие данные ради скорости.\\
    Осознанное: измерено, выбрана ступень, в ключе всё, решён срок.}
\end{frame}

% ============================================================ 2. LOGGING ==
\section{2. Логирование}

% 24
\begin{frame}[plain]
  \sectionopener{2}{Логирование}{Лог пишут не на всякий случай, а чтобы
    ответить на вопрос в \hi{три часа ночи}.}
\end{frame}

% 25
\begin{frame}[plain]
  \vote{Три часа ночи. Пользователь говорит, что его заказ пропал. Какой
    вопрос вы задаёте логам?}
\end{frame}

% 26
\begin{frame}{Четыре вопроса}
  \slidetop
  \begin{center}
    \begin{tikzpicture}[font=\small,
      tb/.style={box,minimum width=3.05cm,minimum height=1.9cm}]
      \node[tb] (a) at (0,0)     {\textbf{что пришло}\\[0.3em]\footnotesize запрос, от кого,\\\footnotesize параметры};
      \node[tb,draw=accent,line width=1.0pt] (b) at (3.45,0)   {\textbf{что решили}\\[0.3em]\footnotesize ветка логики,\\\footnotesize что нашли};
      \node[tb] (c) at (6.9,0)   {\textbf{чем ответили}\\[0.3em]\footnotesize статус,\\\footnotesize что отдали};
      \node[tb] (d) at (10.35,0)  {\textbf{сколько заняло}\\[0.3em]\footnotesize где потратили\\\footnotesize время};
    \end{tikzpicture}
  \end{center}

  \slidegap
  \claim{Каждая строка лога отвечает на \hi{один из них} --- иначе её не должно
    быть.}

  \aside{У шума есть цена: система логов берёт деньги за объём, а в три ночи вы
    ищете четыре строки среди тысяч.}
\end{frame}

% 27
\begin{frame}{Решения, а не действия}
  \slidetop
  \versus{Решение}{%
    \begin{cbul}{accent}
      \item «Заказ 42 не найден для пользователя 7»
      \item отвечает на «что решили»
      \item объясняет, \hi{почему} сервис так сделал
    \end{cbul}}{Действие}{%
    \begin{cbul}{alt}
      \item «Вошёл в метод X»
      \item не отвечает ни на один вопрос
      \item стоит денег и внимания
    \end{cbul}}

  \slidegap
  \hrulethin{0.42}
  \rulegap
  \aside{Не найден, не прошёл проверку, выбрана такая-то ветка, внешний сервис
    ответил вот что.}
\end{frame}

% 28
\begin{frame}[fragile]{Категория}
  \slidetop
  \claim{Тип логгера --- это \hi{категория}: по ней фильтруют.}

\begin{lstlisting}[style=csharp,basicstyle=\codesm,escapechar=|]
namespace OrdersService.Orders;

|\mk{l1}|public sealed class OrderService(ILogger<OrderService> logger)
\end{lstlisting}
  \begin{tikzpicture}[remember picture,overlay,every node/.style={rawlab}]
    \node[text=accent] at ([xshift=10.6cm]l1) {через конструктор};
  \end{tikzpicture}

\begin{lstlisting}[style=term,basicstyle=\termsm]
info: OrdersService.Orders.OrderService[0]
      Order 42 created for user 7 in 1.9518 ms
\end{lstlisting}

  \rulegap
  \aside{Имя класса целиком, с пространством имён, --- и переименуется вместе с
    классом. Статическое поле логгера --- \lo{спрятанная} зависимость: в тесте
    не подменить.}
\end{frame}

% 29
\begin{frame}{Уровни}
  \slidetop
  {\renewcommand{\arraystretch}{1.25}
  \begin{tabular}{@{}L{3.2cm}L{10.0cm}@{}}
    \texttt{Trace}            & всё подряд, для отладки конкретной проблемы \\
    \texttt{Debug}            & подробности, полезные разработчику \\
    \hi{\texttt{Information}} & обычный ход дел: те самые четыре вопроса \\
    \texttt{Warning}          & сервис справился, но не так, как ожидал --- их
      \hi{считают} \\
    \lo{\texttt{Error}}       & операция не удалась --- на него ставится оповещение \\
    \texttt{Critical}         & сервис не может продолжать \\
  \end{tabular}}

  \slidegap
  \aside{Уровень --- это порог: всё ниже отбрасывается до форматирования. Строка
    не стоит того, чтобы разбудить дежурного, --- это \lo{не} \texttt{Error}.}
\end{frame}

% 30
\begin{frame}[fragile]{Уровни --- в конфигурации}
  \slidetop
\begin{lstlisting}[style=json,basicstyle=\codesm,escapechar=|]
{
  "Logging": {
    "LogLevel": {
|\mk{c1}|      "Default": "Information",
|\mk{c2}|      "Microsoft.AspNetCore": "Warning",
|\mk{c3}|      "OrdersService.Orders": "Debug"
    }
  }
}
\end{lstlisting}
  \begin{tikzpicture}[remember picture,overlay,every node/.style={rawlab}]
    \node at ([xshift=8.2cm]c1) {умолчание};
    \node at ([xshift=8.2cm]c2) {категория фреймворка тише};
    \node[text=accent] at ([xshift=8.2cm]c3) {своя категория громче};
  \end{tikzpicture}

\begin{lstlisting}[style=bash,basicstyle=\codesm]
Logging__LogLevel__Default=Debug dotnet run
\end{lstlisting}

  \rulegap
  \aside{Тот же ключ переменной окружения, через двойное подчёркивание: отладка
    на час --- \hi{без пересборки} и выкладки.}
\end{frame}

% 31
\begin{frame}[fragile]{Было и стало}
  \slidetop
\begin{lstlisting}[style=csharp,basicstyle=\codesm,escapechar=|]
|\mk{s1}|logger.LogInformation(
    $"Order {order.Id} created for user {userId} in {ms} ms");

|\mk{s2}|logger.LogInformation(
    "Order {OrderId} created for user {UserId} in {ElapsedMs} ms",
|\mk{s3}|    order.Id, userId, ms);
\end{lstlisting}
  \begin{tikzpicture}[remember picture,overlay,every node/.style={rawlab}]
    \node[text=alt] at ([xshift=10.4cm]s1) {было: одна строка};
    \node[text=accent] at ([xshift=10.4cm]s2) {стало: шаблон с именами};
    \node at ([xshift=10.4cm]s3) {значения отдельно};
  \end{tikzpicture}

  \slidegap
  \versus{Шаблон}{%
    \begin{cbul}{accent}
      \item уровень выключен --- не делает ничего
      \item в записи поля: \texttt{OrderId}, \texttt{UserId}
    \end{cbul}}{Интерполяция}{%
    \begin{cbul}{alt}
      \item строка собрана до вызова, всегда
      \item в записи текст: только читать глазами
    \end{cbul}}
\end{frame}

% 32
\begin{frame}[fragile]{Запись --- это данные}
  \slidetop
  \claim{Лог перестаёт быть текстом: его можно \hi{искать и считать}.}

\begin{lstlisting}[style=term,basicstyle=\termsm]
$ jq -c 'select(.State.UserId == "7") | .State | del(."{OriginalFormat}")' app.log
{"OrderId":42,"UserId":"7","ElapsedMs":2.9927}
{"OrderId":42,"UserId":"7","ElapsedMs":2.9927}
$ jq '.State.ElapsedMs | numbers' app.log | jq -s 'add / length'
2.9927
\end{lstlisting}

  \rulegap
  \begin{cbul}{accent}
    \item «все записи, где \texttt{UserId} равен 7» --- запрос по полю
    \item «среднее \texttt{ElapsedMs} за час» --- агрегация по полю
    \item \texttt{OrderId} и \texttt{OrderID} --- \lo{два разных} поля: имена ---
      часть контракта
  \end{cbul}
\end{frame}

% 33
\begin{frame}[fragile]{Генератор кода}
  \slidetop
\begin{lstlisting}[style=csharp,basicstyle=\codesm,escapechar=|]
static partial class Log
{
|\mk{m1}|    [LoggerMessage(Level = LogLevel.Information,
        Message = "Order {OrderId} created for user {UserId} in {ElapsedMs} ms")]
|\mk{m2}|    public static partial void OrderCreated(
        this ILogger logger, int orderId, string userId, double elapsedMs);
}

|\mk{m3}|logger.OrderCreated(order.Id, userId, ms);
\end{lstlisting}
  \begin{tikzpicture}[remember picture,overlay,every node/.style={rawlab}]
    \node at ([xshift=9.2cm]m1) {объявлено один раз};
    \node[text=accent] at ([xshift=9.2cm]m2) {тело пишет компилятор};
    \node at ([xshift=9.2cm]m3) {вызов --- одно слово};
  \end{tikzpicture}

  \slidegap
  \begin{cbul}{accent}
    \item проверка уровня и форматирование без промежуточных объектов
    \item все сообщения в одном месте --- перечень того, что сервис \hi{умеет сказать}
    \item у сообщения может быть номер события: ищут по номеру, а не по фразе
  \end{cbul}
\end{frame}

% 34
\begin{frame}{Что в логе и чего в нём нет}
  \slidetop
  \versus{Логировать}{%
    \begin{cbul}{accent}
      \item запрос: метод, путь, статус, длительность
      \item решения
      \item исключение --- объектом, один раз, на границе
      \item старт: версия, окружение, адреса
    \end{cbul}}{Никогда}{%
    \begin{cbul}{alt}
      \item секреты, токены, пароли
      \item заголовок авторизации
      \item персональные данные: имя, почта, телефон
      \item тело запроса целиком
    \end{cbul}}

  \slidegap
  \hrulethin{0.42}
  \rulegap
  \aside{Идентификатор пользователя можно: это число. Имя нельзя: лог год
    хранится у третьей стороны.}
\end{frame}

% 35
\begin{frame}[plain]
  \vote{На уровне отладки логируется объект пользователя целиком. В проде
    отладка выключена. Всё безопасно?}
\end{frame}

% 36
\begin{frame}{Десять строк, один запрос}
  \slidetop
  {\renewcommand{\arraystretch}{1.2}\small
  \begin{tabular}{@{}L{2.4cm}L{3.6cm}L{6.9cm}@{}}
    \colhead{muted}{Trace Id} & \colhead{muted}{категория} & \colhead{muted}{сообщение} \\
    \texttt{9c1e\ldots}      & middleware        & GET /orders/17 --- 200 за 12 мс \\
    \hi{\texttt{ada2\ldots}} & \hi{обработчик}   & \hi{заказ 42 запрошен} \\
    \texttt{51b7\ldots}      & клиент курсов     & повтор 1 из 3 \\
    \hi{\texttt{ada2\ldots}} & \hi{сервис}       & \hi{заказ 42 не найден для пользователя 7} \\
    \texttt{9c1e\ldots}      & сервис            & заказ 17 отдан из кэша \\
    \hi{\texttt{ada2\ldots}} & \hi{обработчик ошибок} & \hi{GET /orders/42 --- 404 за 3 мс} \\
    \texttt{51b7\ldots}      & клиент курсов     & ответ за 240 мс \\
  \end{tabular}}

  \slidegap
  \aside{Сто запросов одновременно --- строки вперемешку. Фильтр по одному полю
    --- и вот \hi{история одного запроса}. Клиент присылает этот идентификатор
    из ответа об ошибке.}
\end{frame}

% 37
\begin{frame}[fragile]{Область}
  \slidetop
\begin{lstlisting}[style=csharp,basicstyle=\codesm,escapechar=|]
|\mk{b1}|using (logger.BeginScope(
    new Dictionary<string, object> { ["OrderId"] = orderId }))
{
|\mk{b2}|    logger.OrderCreated(orderId, userId, ms);
}
\end{lstlisting}
  \begin{tikzpicture}[remember picture,overlay,every node/.style={rawlab}]
    \node[text=accent] at ([xshift=9.6cm]b1) {поле для всех строк внутри};
    \node at ([xshift=9.6cm]b2) {любой класс, любой уровень};
  \end{tikzpicture}

\begin{lstlisting}[style=term,basicstyle=\termsm,escapechar=`]
$ jq -c 'select(.EventId == 0) | .Scopes[] | del(.Message, .ParentId)' app.log
`\mk{b3}`{"SpanId":"abc15d1efe7e7268","TraceId":"ada2952ebd54d8219ea277f76698308f"}
{"ConnectionId":"0HNQ0G6CDUHAJ"}
`\mk{b4}`{"RequestId":"0HNQ0G6CDUHAJ:00000001","RequestPath":"/orders/42"}
`\mk{b5}`{"OrderId":42}
\end{lstlisting}
  \begin{tikzpicture}[remember picture,overlay,every node/.style={rawlab}]
    \draw[rawbrace] ([xshift=10.35cm,yshift=0.35em]b3) -- ([xshift=10.35cm,yshift=-0.3em]b4);
    \node at ([xshift=10.6cm]$(b3)!0.5!(b4)$) {открыл фреймворк};
    \node[text=accent] at ([xshift=10.6cm]b5) {открыли вы};
  \end{tikzpicture}

  \aside{Область переживает \texttt{await}: она течёт с асинхронной цепочкой, а
    не с потоком. Область на запрос и на операцию --- не на итерацию цикла.}
\end{frame}

% 38
\begin{frame}{Куда всё это пишется}
  \slidetop
  \claim{Логи --- это \hi{поток событий}, а не файлы.}

  \begin{center}
    \begin{tikzpicture}[font=\footnotesize,
      tb/.style={box,minimum width=3.2cm,minimum height=1.1cm}]
      \node[tb,draw=accent,line width=1.0pt] (s) at (0,0) {\textbf{сервис}\\стандартный вывод};
      \node[tb] (e) at (5.1,0) {\textbf{среда}\\оркестратор, агент};
      \node[tb] (l) at (10.2,0) {\textbf{система логов}\\индекс по полям};
      \draw[arra] (s) -- node[above,font=\scriptsize,text=accent] {JSON, строка} (e);
      \draw[arr] (e) -- (l);
    \end{tikzpicture}
  \end{center}

  \slidegap
  \begin{cbul}{accent}
    \item сервис не знает адреса системы логов и не держит с ней соединение
    \item файл в контейнере исчезнет вместе с ним; ротация --- \lo{не ваша} забота
    \item время --- всемирное; аудит --- отдельный поток со своим хранением
  \end{cbul}
\end{frame}

% 39
\begin{frame}[plain]
  \keyline{Лог отвечает на четыре вопроса.\\
    Он структурный, связан областями, и в нём нет ничего, что нельзя показать
    незнакомцу.}
\end{frame}

% ====================================================== 3. OBSERVABILITY ==
\section{3. Наблюдаемость}

% 40
\begin{frame}[plain]
  \sectionopener{3}{Наблюдаемость}{\hi{Логи} --- что случилось, \hi{метрики} ---
    сколько и как часто, \hi{трейсы} --- где потратили время.}
\end{frame}

% 41
\begin{frame}[plain]
  \vote{Запрос занял три секунды. Где именно?}
\end{frame}

% 42
\begin{frame}{Три сигнала}
  \slidetop
  {\renewcommand{\arraystretch}{1.55}
  \begin{tabular}{@{}L{1.9cm}L{3.2cm}L{3.6cm}L{3.3cm}@{}}
    & \colhead{accent}{вопрос} & \colhead{accent}{что это} & \colhead{alt}{цена} \\
    \textbf{логи}    & что случилось       & события с контекстом
      & дорогие в хранении \\
    \textbf{метрики} & сколько и как часто & числа, накопленные во времени
      & дешёвые: число на интервал \\
    \textbf{трейсы}  & где потратили время & запрос как дерево отрезков
      & объект и экспорт на отрезок \\
  \end{tabular}}

  \slidegap
  \aside{Считать миллионы строк ради одного числа --- медленно, дорого и
    приблизительно. Нужны \hi{все три}: они отвечают на разные вопросы.}
\end{frame}

% 43
\begin{frame}{Водопад}
  \slidetop
  \claim{Вопрос «где» отвечается \hi{картинкой}, а не арифметикой по строкам.}

  \begin{center}
    \begin{tikzpicture}[font=\scriptsize,x=3.1cm,y=0.62cm]
      \foreach \x in {0,1,2,3} {
        \draw[hairline,line width=0.4pt] (\x,0.6) -- (\x,-3.7);
        \node[text=muted,above] at (\x,0.6) {\x\ с};
      }
      \fill[codebg,draw=hairline,line width=0.6pt] (0,0.25) rectangle (3,-0.25);
      \node[anchor=east] at (-0.05,0) {\texttt{GET /orders/42}};
      \fill[codebg,draw=hairline,line width=0.6pt] (0.03,-0.75) rectangle (2.97,-1.25);
      \node[anchor=east] at (-0.05,-1) {обработчик};
      \fill[accent] (0.06,-1.75) rectangle (0.26,-2.25);
      \node[anchor=east] at (-0.05,-2) {\texttt{SELECT orders}};
      \node[anchor=west,text=accent] at (0.28,-2) {0,2 с};
      \fill[alt] (0.3,-2.75) rectangle (2.8,-3.25);
      \node[anchor=east] at (-0.05,-3) {\texttt{GET rates}};
      \node[anchor=west,text=alt] at (2.82,-3) {\textbf{2,5 с}};
    \end{tikzpicture}
  \end{center}

  \aside{Отрезок --- span, дерево --- trace. У отрезков стандартные атрибуты:
    метод, путь, статус, адрес базы --- одинаковые для всех языков.}
\end{frame}

% 44
\begin{frame}{Мониторинг и наблюдаемость}
  \slidetop
  \versus{Наблюдаемость}{%
    \begin{cbul}{accent}
      \item вопросы, которые вы \hi{не задавали}
      \item данные с контекстом, резать в любом направлении
      \item переходы между сигналами по \texttt{TraceId}
      \item от симптома --- к причине
    \end{cbul}}{Мониторинг}{%
    \begin{cbul}{alt}
      \item вопросы, заданные заранее
      \item график и порог оповещения
      \item ловит известное
      \item часть наблюдаемости
    \end{cbul}}

  \slidegap
  \hrulethin{0.42}
  \rulegap
  \aside{Оповещение $\to$ график $\to$ трейсы за этот момент $\to$ логи по
    \texttt{TraceId}. Причина найдена за пять минут, без «добавим логов и
    перевыложим».}
\end{frame}

% 45
\begin{frame}{OpenTelemetry}
  \slidetop
  \claim{Инструментируем \hi{один раз}, экспортируем куда угодно.}

  \begin{center}
    \begin{tikzpicture}[font=\footnotesize,
      tb/.style={box,minimum width=2.7cm,minimum height=0.8cm}]
      \node[tb] (c1) at (0,0.95)  {ваш код};
      \node[tb] (c2) at (0,0)     {ASP.NET Core, HttpClient};
      \node[tb] (c3) at (0,-0.95) {драйвер базы, рантайм};
      \node[box,draw=accent,line width=1.0pt,minimum width=3.0cm,minimum height=2.7cm]
        (o) at (4.3,0) {\textbf{OpenTelemetry}\\[0.3em]модель\\инструментирование\\экспорт};
      \node[tb,minimum width=2.2cm] (k) at (8.0,0) {коллектор};
      \node[tb,minimum width=2.0cm] (s1) at (11.0,0.95)  {система A};
      \node[tb,minimum width=2.0cm] (s2) at (11.0,0)     {система B};
      \node[tb,minimum width=2.0cm] (s3) at (11.0,-0.95) {система C};
      \foreach \n in {c1,c2,c3} \draw[arr] (\n.east) -- (\n.east -| o.west);
      \draw[arra] (o) -- node[above,font=\scriptsize,text=accent] {OTLP} (k);
      \foreach \n in {s1,s2,s3} \draw[arr] (k.east) -- (\n.west);
    \end{tikzpicture}
  \end{center}

  \slidegap
  \aside{Куда уезжают сигналы, решает конфигурация: адрес получателя. Сменили
    систему хранения --- поменяли адрес, код \hi{не тронут}.}
\end{frame}

% 46
\begin{frame}[fragile]{Регистрация}
  \vspace{0.3em}
\begin{lstlisting}[style=csharp,basicstyle=\codexs,escapechar=|]
builder.Services.AddOpenTelemetry()
|\mk{t1}|    .ConfigureResource(r => r.AddService("orders-service"))
    .WithTracing(tracing => tracing
|\mk{t2}|        .AddAspNetCoreInstrumentation()
|\mk{t3}|        .AddHttpClientInstrumentation()
|\mk{t4}|        .AddSource("OrdersService")
|\mk{t5}|        .AddOtlpExporter())
    .WithMetrics(metrics => metrics
        .AddAspNetCoreInstrumentation()
        .AddHttpClientInstrumentation()
|\mk{t6}|        .AddRuntimeInstrumentation()
|\mk{t7}|        .AddMeter("OrdersService")
        .AddOtlpExporter());

|\mk{t8}|builder.Logging.AddOpenTelemetry(
    logging => logging.AddOtlpExporter());
\end{lstlisting}
  \begin{tikzpicture}[remember picture,overlay,every node/.style={rawlab}]
    \node at ([xshift=9.6cm]t1) {имя сервиса};
    \node at ([xshift=9.6cm]t2) {входящие запросы};
    \node at ([xshift=9.6cm]t3) {исходящие вызовы};
    \node at ([xshift=9.6cm]t4) {свои отрезки};
    \node[text=accent] at ([xshift=9.6cm]t5) {адрес --- из окружения};
    \node at ([xshift=9.6cm]t6) {рантайм};
    \node at ([xshift=9.6cm]t7) {свои метрики};
    \node[text=accent] at ([xshift=9.6cm]t8) {и логи туда же};
  \end{tikzpicture}

  \aside{Ни строки в вашем коде: запрос --- отрезок, вызов клиента ---
    \hi{дочерний} отрезок.}
\end{frame}

% 47
\begin{frame}{Цена}
  \slidetop
  \versus{Трейсы сэмплируют}{%
    \begin{cbul}{accent}
      \item отрезок --- объект и экспорт по сети
      \item записывают долю: скажем, 10\,\%
      \item решение в начале --- на всю цепочку
      \item в конце: все медленные и с ошибками
    \end{cbul}}{Метрики: кардинальность}{%
    \begin{cbul}{alt}
      \item сочетание атрибутов --- отдельный ряд
      \item тип оплаты --- три ряда
      \item идентификатор пользователя --- миллион
      \item в атрибутах --- десятки значений
    \end{cbul}}

  \slidegap
  \hrulethin{0.42}
  \rulegap
  \aside{Метрики не сэмплируют никогда. Идентификаторы пользователей и заказов
    живут в логах и трейсах, а \lo{не} в метриках.}
\end{frame}

% 48
\begin{frame}[fragile]{Свои отрезки и метрики}
  \vspace{0.3em}
\begin{lstlisting}[style=csharp,basicstyle=\codexxs,escapechar=|]
static class Telemetry
{
|\mk{y1}|    public static readonly ActivitySource Source =
        new("OrdersService");
    public static readonly Meter Meter = new("OrdersService");
    public static readonly Counter<long> OrdersCreated =
        Meter.CreateCounter<long>("orders.created");
}

|\mk{y2}|using (var activity =
    Telemetry.Source.StartActivity("pricing.calculate"))
{
|\mk{y3}|    activity?.SetTag("order.id", orderId);
    // ...
|\mk{y4}|    Telemetry.OrdersCreated.Add(1,
        new KeyValuePair<string, object?>("payment", "card"));
}
\end{lstlisting}
  \begin{tikzpicture}[remember picture,overlay,every node/.style={rawlab}]
    \node at ([xshift=9.6cm]y1) {один на сервис};
    \node[text=accent] at ([xshift=9.6cm]y2) {отрезок вокруг логики};
    \node at ([xshift=9.6cm]y3) {\texttt{?}: отрезка может не быть};
    \node at ([xshift=9.6cm]y4) {счётчик с атрибутом};
  \end{tikzpicture}

  \aside{Только типы платформы: OpenTelemetry подхватывает их \hi{по именам}.}
\end{frame}

% 49
\begin{frame}[plain]
  \vote{Заказы вызвали курсы, курсы --- справочники. Медленно. Как понять, что
    три отрезка в трёх сервисах --- один запрос пользователя?}
\end{frame}

% 50
\begin{frame}{Распространение контекста}
  \slidetop
  \begin{center}
    \begin{tikzpicture}[font=\footnotesize,
      tb/.style={box,minimum width=2.6cm,minimum height=1.0cm}]
      \node[tb] (a) at (0,0)   {\textbf{заказы}};
      \node[tb] (b) at (5.2,0) {\textbf{курсы}};
      \node[tb] (c) at (10.4,0) {\textbf{справочники}};
      \draw[arra] (a) -- node[above,font=\scriptsize\ttfamily,text=accent] {traceparent} (b);
      \draw[arra] (b) -- node[above,font=\scriptsize\ttfamily,text=accent] {traceparent} (c);
      \node[font=\small\ttfamily,below=1.0em] at (b.south)
        {traceparent: 00-\hi{ada2952ebd54d8219ea277f76698308f}-abc15d1efe7e7268-01};
      \node[font=\scriptsize,text=accent,below=2.9em] at (b.south)
        {идентификатор трейса --- один на всю цепочку};
    \end{tikzpicture}
  \end{center}

  \slidegap
  \begin{cbul}{accent}
    \item клиент из фабрики добавляет заголовок \hi{сам}; сервер продолжает тот же трейс
    \item очередь: идентификатор кладут в сообщение и восстанавливают
    \item багаж виден всем по пути --- \lo{ничего} чувствительного
  \end{cbul}
\end{frame}

% 51
\begin{frame}{Какой сигнал}
  \slidetop
  {\renewcommand{\arraystretch}{1.5}
  \begin{tabular}{@{}L{2.6cm}L{3.8cm}L{6.4cm}@{}}
    \hi{\textbf{метрика}} & «что-то не так»  & счётчик ошибок: дёшево, мгновенно,
      оповещение \\
    \hi{\textbf{лог}}     & «что именно»     & какие ошибки, когда счётчик подскочил \\
    \hi{\textbf{трейс}}   & «где»            & на каком отрезке и в каком сервисе \\
  \end{tabular}}

  \slidegap
  \hrulethin{0.42}
  \rulegap
  \eyebrow{четыре золотых сигнала}\par
  \vspace{0.5em}
  {\large \hi{задержка} \quad \hi{трафик} \quad \hi{ошибки} \quad \hi{насыщение}}

  \vspace{0.5em}
  \aside{Четыре графика на одной панели, одинаковой для всех сервисов.}
\end{frame}

% 52
\begin{frame}[plain]
  \keyline{Три сигнала, один идентификатор, одна разметка.\\
    Наблюдаемость --- когда от любого симптома можно дойти до причины.}
\end{frame}

% ============================================================= 4. ASPIRE ==
\section{4. Оркестрация локального окружения}

% 53
\begin{frame}[plain]
  \sectionopener{4}{Оркестрация локального окружения}{Сервис, база и Redis
    \hi{одной командой} --- и всё это видно в одном окне.}
\end{frame}

% 54
\begin{frame}[plain]
  \vote{Что нужно поднять, чтобы запустить наш сервис локально?}
\end{frame}

% 55
\begin{frame}[fragile]{Хост приложения}
  \slidetop
  \claim{Система описана \hi{кодом}, в отдельном проекте.}

\begin{lstlisting}[style=csharp,basicstyle=\codesm,escapechar=|]
var builder = DistributedApplication.CreateBuilder(args);

|\mk{a1}|var cache = builder.AddRedis("cache");
|\mk{a2}|var database = builder.AddPostgres("postgres")
    .AddDatabase("orders");

builder.AddProject<Projects.OrdersService>("orders")
|\mk{a3}|    .WithReference(cache)
|\mk{a4}|    .WithReference(database);

|\mk{a5}|builder.Build().Run();
\end{lstlisting}
  \begin{tikzpicture}[remember picture,overlay,every node/.style={rawlab}]
    \node at ([xshift=9.4cm]a1) {Redis как ресурс};
    \node at ([xshift=9.4cm]a2) {база как ресурс};
    \draw[rawbrace] ([xshift=9.15cm,yshift=0.35em]a3) -- ([xshift=9.15cm,yshift=-0.3em]a4);
    \node[text=accent] at ([xshift=9.4cm]$(a3)!0.5!(a4)$) {сервис ссылается};
    \node at ([xshift=9.4cm]a5) {одна команда};
  \end{tikzpicture}

  \aside{Сервис не знает портов и адресов: строки подключения \texttt{cache} и
    \texttt{orders} приходят через окружение.}
\end{frame}

% 56
\begin{frame}{Что делает хост}
  \slidetop
  \claim{Всё, что вы делали бы руками в пяти терминалах.}

  \begin{cbul}{accent}
    \item поднимает контейнеры, порты выбирает сам
    \item запускает проекты сервисов как обычные процессы
    \item передаёт каждому сервису \hi{переменные окружения}
    \item следит за состоянием и останавливает всё вместе с собой
  \end{cbul}

  \slidegap
  {\renewcommand{\arraystretch}{1.2}\small
  \begin{tabular}{@{}L{6.2cm}L{7.0cm}@{}}
    \texttt{ConnectionStrings\_\_cache}        & строка подключения к ресурсу \\
    \texttt{OTEL\_EXPORTER\_OTLP\_ENDPOINT}    & адрес панели для экспорта сигналов \\
    \texttt{OTEL\_SERVICE\_NAME}               & имя сервиса \\
  \end{tabular}}
\end{frame}

% 57
\begin{frame}[fragile]{Умолчания и интеграции}
  \slidetop
\begin{lstlisting}[style=csharp,basicstyle=\codesm,escapechar=|]
|\mk{d1}|builder.AddServiceDefaults();
|\mk{d2}|builder.AddRedisDistributedCache("cache");
\end{lstlisting}
  \begin{tikzpicture}[remember picture,overlay,every node/.style={rawlab}]
    \node at ([xshift=8.4cm]d1) {проект умолчаний};
    \node[text=accent] at ([xshift=8.4cm]d2) {интеграция: имя как в хосте};
  \end{tikzpicture}

  \slidegap
  \versus{Проект умолчаний}{%
    \begin{cbul}{accent}
      \item OpenTelemetry: трейсы, метрики, логи
      \item проверки здоровья
      \item устойчивость клиентов
    \end{cbul}}{Пакет интеграции}{%
    \begin{cbul}{alt}
      \item клиент приходит из контейнера
      \item строка подключения настроена
      \item проверка здоровья и разметка трейсов
    \end{cbul}}
\end{frame}

% 58
\begin{frame}{Панель}
  \slidetop
  \claim{Трейсы своего сервиса --- на \hi{второй минуте}, без инфраструктуры.}

  {\renewcommand{\arraystretch}{1.4}
  \begin{tabular}{@{}L{3.0cm}L{10.2cm}@{}}
    \hi{\textbf{ресурсы}} & все ли работают; консольный вывод с ошибкой старта \\
    \hi{\textbf{логи}}    & структурные, с полями, фильтр по уровню и тексту \\
    \hi{\textbf{трейсы}}  & водопад запроса с отрезками Redis и базы \\
    \hi{\textbf{метрики}} & длительность запросов, память, пул потоков, ваш счётчик \\
  \end{tabular}}

  \slidegap
  \aside{Aspire --- про локальную разработку, \lo{не про прод}: там нужна
    настоящая система хранения сигналов. И нужен запущенный движок контейнеров.}
\end{frame}

% 59
\begin{frame}[plain]
  \keyline{Система описана кодом и поднята одной командой.\\
    Наблюдаемость --- с первого дня, а не после первого инцидента.}
\end{frame}

\end{document}
